% !TeX root = ../../../金融经济学_王江_讲义及习题.tex
% 同一三证券模型的示意：预算平面坐标(t,u)，均值=.4+.4t，
% 方差=.36+.36t^2+.64u^2；前沿u=0，MVP=(0,0)，p=(.8,.8)。
% 混合权重(t,u)=(.8 alpha,.8 alpha)，方差=.36+.64 alpha^2。
\begin{tikzpicture}[>=Stealth,font=\small,
 frontier/.style={blue!70!black,very thick},
 mixing/.style={teal!65!black,very thick},
 guide/.style={gray,densely dotted}]
 \begin{scope}
  \draw[->] (0,0)--(4.35,0) node[right] {$\sigma$};
  \draw[->] (0,0)--(0,4.05) node[above] {$\bar r$};
  \draw[frontier] plot[domain=-1.2:1.8,samples=81,variable=\t]
   ({2.4*sqrt(.36+.36*\t*\t)},{.6+2.6*(.4+.4*\t)});
  \draw[mixing] plot[domain=-1.5:1.85,samples=81,variable=\a]
   ({2.4*sqrt(.36+.64*\a*\a)},{.6+2.6*(.4+.32*\a)});
  \draw[guide] (0,1.64)--(1.44,1.64)--(1.44,0);
  \fill[blue!70!black] ({2.4*sqrt(.72)},.6) circle(2pt)
   node[below left=2pt] {$z_0$};
  \fill[blue!70!black] ({2.4*sqrt(1.17)},3.2) circle(2pt)
   node[above left=2pt] {$z_1$};
  \fill (1.44,1.64) circle(2.2pt) node[left=4pt] {$\mathrm{MVP}$};
  \fill[teal!65!black] (2.4,2.472) circle(2.2pt) node[right=4pt] {$p$};
  \draw[frontier] (.25,3.85)--(.65,3.85) node[right,black] {原市场前沿};
  \draw[mixing] (2.45,1.15)--(2.85,1.15) node[right,black] {$p$ 与 MVP 混合};
  \node at (2.1,-.55) {(a) 均值--标准差平面};
 \end{scope}
 \begin{scope}[xshift=6.5cm,yshift=.25cm]
  \filldraw[fill=blue!5,draw=gray!65]
   (0,.3)--(3.3,.3)--(5.1,3.5)--(1.8,3.5)--cycle;
  \node[above] at (2.7,3.57) {三证券的预算平面 $\iota^{\mathsf T}z=1$};
  \draw[frontier,<->] (.95,.9)--(4.1,3);
  \draw[mixing,<->] (.5,1.7)--(4.6,1.928);
  \fill[blue!70!black] (1.4,1.2) circle(2pt) node[below right=2pt] {$z_0$};
  \fill[blue!70!black] (3.65,2.7) circle(2pt) node[above left=2pt] {$z_1$};
  \fill (2.3,1.8) circle(2.2pt) node[above left=3pt] {$z_{\mathrm{MVP}}$};
  \fill[teal!65!black] (3.74,1.88) circle(2.2pt) node[below=4pt] {$z_p$};
  \node[blue!70!black,rotate=34] at (3.3,2.2) {前沿直线};
  \node[teal!65!black] at (2.6,.43) {$z_\alpha=(1-\alpha)z_{\mathrm{MVP}}+\alpha z_p$};
  \node at (2.55,-.8) {(b) 权重空间中的预算平面};
 \end{scope}
\end{tikzpicture}
